\chapter{线性代数}\label{ch:b2:linalg}

本科一年级分册中的线性代数是在 $\R$ 或 $\C$ 上
讨论有限维空间的，并对一般行列式仅作了承认。
本章将三个限制全部推广：理论建立在任意域
$K$ 上，空间与其\emph{对偶空间}之间的关系得到系统展开
（对偶基、零化子、转置），行列式则最终由
交错多重线性形式及\cref{ch:b2:structures}中的符号\emph{构造}而来，
兑现了本科一年级分册中的
全部承认。

全文中，$K$ 是一个域（$\Q$、$\R$、$\C$ 或 $\Z/p\Z$ --- 理论
并不依赖于具体选取），除非另有说明，空间均为
$K$ 上的有限维空间。本科一年级分册的结果（基、维数、
秩--零度定理、矩阵）原样移用：其证明
除域的公理外从未用到其他任何东西。

\section{对偶空间}

\begin{definition}[对偶空间，对偶基]\label{def:b2:linalg:dual}
设 $E$ 的\emph{对偶空间}\index{对偶空间}为 $E^* = \mathcal{L}(E,
K)$，即全体线性型构成的空间。若 $\mathcal{B} = (e_1, \dots, e_n)$
为 $E$ 的一组基，其\emph{坐标线性型} $e_1^*, \dots, e_n^*$
由 $e_i^*(e_j) = \delta_{ij}$ 定义（克罗内克符号：$1$ 当 $i = j$ 时，
否则为 $0$），构成\emph{对偶基}\index{对偶基}
$\mathcal{B}^*$，它是 $E^*$ 的一组基；特别地，$\dim E^* = \dim E$，且有
\[
x = \sum_{i=1}^{n} e_i^*(x)\, e_i \quad (x \in E),
\qquad
\varphi = \sum_{i=1}^{n} \varphi(e_i)\, e_i^* \quad (\varphi \in
E^*).
\]
\end{definition}

\begin{proof}[证明 $\mathcal{B}^*$ 是一组基]
线性无关：将零组合 $\sum \lambda_i e_i^* = 0$ 作用于
$e_j$ 得 $\lambda_j = 0$。生成：对 $\varphi \in E^*$，
线性型 $\varphi - \sum_i \varphi(e_i) e_i^*$ 在每个 $e_j$ 上取值为零，
故其为零（在一组基上取零的线性映射必为零）。两个展开式
不过是同一计算正着读的结果。
\end{proof}

\begin{example}\label{ex:b2:linalg:dualexamples}
在 $K_n[X]$ 中取基 $(1, X, \dots, X^n)$，其对偶基为 $P
\mapsto \frac{P^{(k)}(0)}{k!}$（即泰勒系数）。对偶空间的另一组基由
求值 $P \mapsto P(x_i)$（取 $n + 1$ 个互异点）给出，
它在 $K_n[X]$ 中的“前对偶”基恰为拉格朗日多项式族
$L_i$（见本科一年级分册），因为 $L_i(x_j) =
\delta_{ij}$。插值\emph{即}对偶。
\end{example}

\begin{method}[对偶基与前对偶基的实际应用]\label{met:b2:linalg:dualbasis}
要将线性型 $\varphi$ 在基 $(e_i)$ 上展开（$E$ 的一组基）：
其坐标就是\emph{值} $\varphi(e_i)$ --- 无需
解任何方程组。求 $(u_j)$（$E$ 的基），使其对偶基恰为给定的
基 $(\varphi_1, \dots, \varphi_n)$（$E^*$ 的一组基，
即\emph{前对偶基}）：解 $n$ 个线性方程组
\[
\varphi_i(u_j) = \delta_{ij} \qquad (1 \leq i \leq n),
\]
逐列求出 $u_j$；用矩阵语言说，若 $M$ 的各行
列出 $\varphi_i$ 在 $E^*$ 的某已知基下的系数，则
$M^{-1}$ 的各列即为 $u_j$。前对偶基的存在性与
唯一性在本章的周末问题中证明；
计算始终归结为这一求逆。
\end{method}

\begin{example}[$\R^2$ 的一组对偶基（完整计算）]\label{ex:b2:linalg:dualbasis2d}
对于 $b_1 = (1, 1)$、$b_2 = (1, -1)$ 构成的 $\R^2$ 的基，其对偶基
$(b_1^*, b_2^*)$ 应满足 $b_i^*(b_j) =
\delta_{ij}$。设 $b_1^*(x, y) = \alpha x + \beta y$，
由条件 $\alpha + \beta = 1$ 与 $\alpha - \beta = 0$ 得
\[
b_1^*(x, y) = \frac{x + y}{2},
\qquad\text{同理}\qquad
b_2^*(x, y) = \frac{x - y}{2} .
\]
验算：$b_1^*$ \emph{不是}把 $e_1^* + e_2^*$ 简单求值所得 ---
对偶基依赖于整组基，而非各个向量单独决定，
（将 $b_2$ 换为 $(0, 1)$ 时，
$b_1^*$ 变为 $x \mapsto x$）。展开式同样成立：
$(x, y) = \frac{x+y}2\,b_1 + \frac{x-y}2\,b_2$，这正是
一对数的奇偶分解 --- 对偶基是坐标提取器，
本例中这组对偶基提取的是对称与
反对称两部分。
\end{example}

\begin{definition}[零化子]\label{def:b2:linalg:annihilator}
设 $F \subseteq E$ 为一子空间，其
\emph{零化子}\index{零化子}为
\[
F^{\circ} = \{\varphi \in E^* : \varphi|_F = 0\},
\]
它是 $E^*$ 的子空间。
\end{definition}

\begin{theorem}[零化子的维数]\label{thm:b2:linalg:annihilator}
$\dim F^{\circ} = \dim E - \dim F$。此外，映射 $F \mapsto F^\circ$
反转包含关系，并且 $F$ 可由其零化子恢复：
\[
F = \{x \in E : \forall\varphi \in F^\circ,\ \varphi(x) = 0\}.
\]
由此，任一 $p$ 维子空间在 $n$ 维空间中都是
$n - p$ 个独立线性方程的解集 ---
反之亦然。
\end{theorem}

\begin{proof}
设 $(e_1, \dots, e_p)$ 为 $F$ 的一组基，将其扩充为
$E$ 的一组基。线性型 $\varphi = \sum \varphi(e_i) e_i^*$ 零化 $F$，当且仅当
其前 $p$ 个系数为零：$F^\circ =
\operatorname{Vect}(e_{p+1}^*, \dots, e_n^*)$，其维数为 $n - p$。
包含关系的反转是显然的。关于恢复，
右端包含 $F$；反之，若 $x \notin F$，则补全 $F$ 的一组基，
用 $x$ 与若干向量扩充；$x$ 在此基下的坐标线性型
零化 $F$ 但不零化 $x$。“方程”的读法是取
$(\varphi_1, \dots, \varphi_{n-p})$ 为 $F^\circ$ 的一组基，则 $F =
\bigcap \ker\varphi_j$，即 $n - p$ 个独立
超平面的交。
\end{proof}

\begin{example}[零化子的两个方向]\label{ex:b2:linalg:annihilatorwork}
设 $F = \operatorname{Vect}\bigl((1, 2, 1),\ (1, 0, -1)\bigr)
\subseteq \R^3$。线性型 $\varphi = a\,e_1^* + b\,e_2^* + c\,e_3^*$
零化 $F$，当且仅当
\[
a + 2b + c = 0
\qquad\text{且}\qquad
a - c = 0 ,
\]
即 $c = a$ 且 $b = -a$：$F^\circ = \R\,(e_1^* - e_2^* +
e_3^*)$，其维数为 $3 - 2 = 1$，
正如 \cref{thm:b2:linalg:annihilator} 所要求。逆向读出：
$F = \{(x, y, z) : x - y + z = 0\}$ --- 这个平面可表示为
张成 $F^\circ$ 的唯一线性型的核。由
张成族到方程组的过程，\emph{正是}计算零化子；
由方程组到参数化的过程，则是计算
前零化子。（验证：两个张成向量都满足 $x - y +
z = 0$。）
\end{example}

\begin{definition}[转置映射]\label{def:b2:linalg:transpose}
对 $u \in \mathcal{L}(E, F)$，称 \emph{转置}\index{转置
映射} $u^{\mathsf T} \in \mathcal{L}(F^*, E^*)$ 为
\[
u^{\mathsf T}(\psi) = \psi \circ u .
\]
且满足 $(v \circ u)^{\mathsf T} = u^{\mathsf T} \circ v^{\mathsf
T}$。在对偶基下，$u^{\mathsf T}$ 的矩阵即
$u$ 的转置矩阵 --- 这\emph{解释}了
本科一年级所学的转置。
\end{definition}

\begin{example}[转置，逐项计算]\label{ex:b2:linalg:transposework}
设 $u \colon \R^2 \to \R^3$ 在标准基下的矩阵为 $A =
\begin{psmallmatrix} 1 & 2\\ 0 & 1\\ 3 & 0\end{psmallmatrix}$。
对 $\psi = b_1f_1^* + b_2f_2^* + b_3f_3^*
\in (\R^3)^*$，计算 $u^{\mathsf T}(\psi) = \psi \circ u$ 在
$\R^2$ 的基上：
\[
(\psi \circ u)(e_1) = \psi(1, 0, 3) = b_1 + 3b_3,
\qquad
(\psi \circ u)(e_2) = \psi(2, 1, 0) = 2b_1 + b_2 .
\]
于是 $u^{\mathsf T}(\psi) = (b_1 + 3b_3)\,e_1^* + (2b_1 +
b_2)\,e_2^*$。在对偶基下，$u^{\mathsf T}$ 的矩阵
为
\[
\begin{pmatrix} 1 & 0 & 3\\ 2 & 1 & 0\end{pmatrix}
= A^{\mathsf T} :
\]
抽象的转置\emph{正是}翻转后的矩阵，无需
再凭信念接受任何计算。注意其机理：
第 $j$ \emph{列}（$A$ 的）变成了
新矩阵的第 $j$ \emph{行}，因为 $\psi \circ u$ 读取 $u$ 的
输出时借助 $\psi$ 的系数。
\end{example}

\begin{proposition}\label{prop:b2:linalg:transposerank}
$\ker u^{\mathsf T} = (\operatorname{im} u)^{\circ}$，且
$\operatorname{im} u^{\mathsf T} = (\ker u)^{\circ}$。因此
$\operatorname{rk}(u^{\mathsf T}) = \operatorname{rk}(u)$：行秩
等于列秩，由结构性论证得到。
\end{proposition}

\begin{proof}
$\psi \in \ker u^{\mathsf T} \iff \psi \circ u = 0 \iff \psi$ 零化
$\operatorname{im} u$，此为第一个等式。对于第二个，$u^{\mathsf
T}(\psi) = \psi \circ u$ 总零化 $\ker u$，故
$\operatorname{im} u^{\mathsf T} \subseteq (\ker u)^\circ$；
维数相等，由秩–零度定理及
\cref{thm:b2:linalg:annihilator}，可得：
\[
\operatorname{rk} u^{\mathsf T} = \dim F^* - \dim\ker u^{\mathsf T}
= \dim F - \bigl(\dim F - \operatorname{rk} u\bigr)
= \operatorname{rk} u
= \dim (\ker u)^{\circ} . \qedhere
\]
\end{proof}

\begin{example}[从两侧读秩]\label{ex:b2:linalg:rankbothsides}
设
\[
A = \begin{pmatrix}
1 & 2 & 0 & 1\\
0 & 1 & 1 & 1\\
1 & 3 & 1 & 2
\end{pmatrix} .
\]
\emph{列秩：}第三行是前两行之和，因此
$\operatorname{rk} A \leq 2$；第 $1$ 列与第 $2$ 列线性无关：
$\operatorname{rk} A = 2$。\emph{转置的核：}解
$A^{\mathsf T}y = 0$ 得 $y \in \R\,(1, 1, -1)$，故 $\ker
A^{\mathsf T}$ 的维数为 $1 = 3 - 2$，恰好是
$(\operatorname{im} A)^\circ$，在 $(\R^3)^*$ 与行向量的
对应下，正如
\cref{prop:b2:linalg:transposerank} 所断言 --- 唯一的
单一关系“行$_3$ = 行$_1$ + 行$_2$”\emph{正是}
列空间的零化子。行秩（$2$ 个线性无关的行）与
列秩相等，并非偶然，而是因为两者都等于
$\operatorname{rk} A = \operatorname{rk} A^{\mathsf T}$。
\end{example}

\begin{example}[对偶性读出求积公式]\label{ex:b2:linalg:quadratureread}
为何像辛普森求积公式（见\cref{exo:b2:linalg:4}）这样的公式存在，
且为何唯一？对偶性在任何计算之前就给出了答案。
在 $E = \R_2[X]$ 上，积分 $P \mapsto \int_0^1 P$ 是
三维对偶空间 $E^*$ 中的一个特定向量；
在 $0$、$\frac12$、$1$ 处的求值泛函构成
$E^*$ 的\emph{基}；因此积分可唯一地在其上展开 --- 此展开
\emph{正是}辛普森公式，系数一并包含在内。
维数统计也能校准预期：在 $\R_3[X]$ 上，
四维的泛函空间一般无法由三个
求值泛函张成，故对三次多项式的精确性并非对偶性所保证；
辛普森公式居然对三次多项式精确积分，这只是
一个额外的对称性（关于 $\frac12$ 的奇次项相消），
需手工验证。含 $n + 1$ 个节点的公式是
积分形式在 $\R_n[X]^*$ 的求值基下的展开：
存在性与唯一性只需一条对偶基定理；仅
额外的次数需要付出计算代价。
\end{example}

\section{多重线性交错形式}

\begin{definition}\label{def:b2:linalg:alternating}
映射 $f \colon E^n \to K$ 称为\emph{$n$-线性}的，当它关于
每个变元都线性，并称为\emph{交错}\index{交错形式}的，当它在
两个自变元相等时恒为零。交错性蕴含
\emph{反对称}：交换两个自变元会改变符号
（展开 $f(\dots, x + y, \dots, x + y, \dots) = 0$）；更一般地，
对 $\sigma \in \mathfrak{S}_n$，
\[
f(x_{\sigma(1)}, \dots, x_{\sigma(n)}) =
\varepsilon(\sigma)\, f(x_1, \dots, x_n),
\]
将 $\sigma$ 分解为对换，由此即得
（见\cref{thm:b2:structures:signature}）。
\end{definition}

\begin{theorem}[行列式基本定理]\label{thm:b2:linalg:detspace}
设 $\dim E = n$，$\mathcal{B} = (e_1, \dots, e_n)$ 为一组基。
全体交错 $n$-线性形式（定义于 $E$ 上）构成的空间维数为 $1$：
每个此类形式都是下式的倍数：
\[
\det{}_{\mathcal{B}}(x_1, \dots, x_n)
= \sum_{\sigma \in \mathfrak{S}_n} \varepsilon(\sigma)
\prod_{i=1}^{n} a_{\sigma(i),\,i},
\qquad
x_j = \sum_{i} a_{ij} e_i ,
\]
而 $\det_{\mathcal{B}}$ 是唯一取值为 $1$ 的形式，
即在基 $\mathcal{B}$ 上的取值。\index{行列式!构造}
\end{theorem}

\begin{proof}
设 $f$ 为交错 $n$-线性形式，将每个自变元
由多重线性在 $\mathcal{B}$ 上展开，
\[
f(x_1, \dots, x_n)
= \sum_{i_1, \dots, i_n} a_{i_1,1}\cdots a_{i_n,n}\,
f(e_{i_1}, \dots, e_{i_n}).
\]
含重复指标的项为零（由交错性）；保留的
元组 $(i_1, \dots, i_n)$ 即单射的那些，亦即 $i_k =
\sigma(k)$，$\sigma$ 为某个置换；由反对称性重排得
$f(e_{\sigma(1)}, \dots, e_{\sigma(n)}) = \varepsilon(\sigma)
f(e_1, \dots, e_n)$。由此可得
\[
f = f(e_1, \dots, e_n) \cdot \det{}_{\mathcal{B}} :
\]
即每个交错形式都是该倍数，前提是
$\det_{\mathcal{B}}$ 本身（即上式的和）\emph{是}
交错 $n$-线性的，且 $1$ 是其在 $\mathcal B$ 上的值。
多重线性显然成立（每个求和项关于每一列都是线性的）。
在 $\mathcal B$ 上的取值：唯一非零项为 $\sigma =
\mathrm{id}$。交错性：设 $x_j = x_k$（$j \neq k$），于是
各坐标列满足 $a_{i j} = a_{i k}$，对所有
$i$。将每个 $\sigma$ 与 $\sigma' = \sigma\circ(j\,k)$ ---
这是 $\mathfrak{S}_n$ 上无不动点的对合。
成对的两项乘积相等：
\[
\prod_i a_{\sigma'(i),\,i}
= a_{\sigma(k),\,j}\; a_{\sigma(j),\,k}
\prod_{i \neq j,k} a_{\sigma(i),\,i}
= a_{\sigma(k),\,k}\; a_{\sigma(j),\,j}
\prod_{i \neq j,k} a_{\sigma(i),\,i}
= \prod_i a_{\sigma(i),\,i},
\]
利用第 $j$ 列与第 $k$ 列相等，而
$\varepsilon(\sigma') = -\varepsilon(\sigma)$。每一对
贡献为零，故和为零。
\end{proof}

\begin{example}[萨鲁斯法则的推导与推翻]\label{ex:b2:linalg:sarrus}
当 $n = 3$ 时，置换公式恰有 $3! = 6$ 项。
按符号列出 $\mathfrak{S}_3$ 的元素 ---
$\mathrm{id}$、$(1\,2\,3)$、$(1\,3\,2)$ 为偶置换；
$(1\,2)$、$(1\,3)$、$(2\,3)$ 为奇置换 --- 得
\[
\det A = a_{11}a_{22}a_{33} + a_{21}a_{32}a_{13} +
a_{31}a_{12}a_{23}
- a_{21}a_{12}a_{33} - a_{31}a_{22}a_{13} -
a_{11}a_{32}a_{23} :
\]
正是中学所教的萨鲁斯“对角线”法则 ---
如今它已是定理，其中神秘的正负号被识别为
置换的符号。推翻之处在于：当 $n = 4$ 时共有 $24$
个置换，其中仅 $8$ 个能被任何画对角线的方案
选中；萨鲁斯法则没有 $4$ 阶版本，
而余子式展开（见\cref{thm:b2:linalg:detrules}~(4)）
接手。数项也是一个警示：置换公式
含 $n!$ 个求和项，故它是\emph{定义}，而非
算法 --- 行化简计算 $\det$ 只需
$O(n^3)$ 次运算即可。
\end{example}

\begin{definition}[行列式]\label{def:b2:linalg:det}
在一组基下的\emph{向量组的行列式}\index{行列式}为
$\det_{\mathcal{B}}(x_1, \dots, x_n)$；\emph{矩阵的
行列式}即 $A$ 各列在标准基下的行列式，
--- 即上述置换公式；\emph{自同态的
行列式} $u$ 为满足下式的标量 $\det u$：
\[
\det{}_{\mathcal{B}}\bigl(u(x_1), \dots, u(x_n)\bigr)
= \det u \cdot \det{}_{\mathcal{B}}(x_1, \dots, x_n)
\quad \text{ 对所有 } x_i
\]
（左边是交错 $n$-线性的，因而是
$\det_\mathcal{B}$ 的倍数，由\cref{thm:b2:linalg:detspace}可知；该因子
不依赖于 $\mathcal{B}$）。
\end{definition}

\begin{theorem}[行列式的运算法则（已证）]\label{thm:b2:linalg:detrules}
\begin{enumerate}
  \item $\det(uv) = \det u\,\det v$; $\;\det(AB) = \det A \det B$.
  \item $u$ 可逆 $\iff \det u \neq 0$；一个向量组为基
        $\iff$ 它在某一组基下的行列式不为零。
  \item $\det(A^{\mathsf T}) = \det A$.
  \item 正如本科一年级分册所述，沿任一行或任一列的余子式展开成立；
        相似矩阵有相同的
        行列式。
\end{enumerate}
\end{theorem}

\begin{proof}
(1) 两次应用定义式：
$\det_{\mathcal B}(uv(x_i)) = \det u \cdot \det_{\mathcal
B}(v(x_i)) = \det u \det v \cdot \det_{\mathcal B}(x_i)$.

(2) 若 $u$ 可逆，$\det u \det u^{-1} = \det \mathrm{id} =
1 \neq 0$。若不可逆，则向量 $u(e_i)$ 线性相关；把其中一个
用其余线性表出并展开，得 $\det_{\mathcal B}(u(e_i)) = 0$
（交错性消去重复方向），因此 $\det u = 0$。基的
判别准则对向量组是同一结论。

(3) 在置换公式中，将每个乘积按 $j =
\sigma(i)$ 换下标，即 $i = \tau(j)$，其中 $\tau = \sigma^{-1}$，
各因子是相同的数，只是次序不同，于是
\[
\prod_{i=1}^{n} a_{\sigma(i),\,i} = \prod_{j=1}^{n}
a_{j,\,\tau(j)} ,
\]
且 $\varepsilon(\tau) = \varepsilon(\sigma)^{-1} =
\varepsilon(\sigma)$（取值为 $\pm1$；$\varepsilon$ 是
同态）。对 $\sigma$ 求和与对
$\tau$ 求和相同（求逆是 $\mathfrak{S}_n$ 的双射）：
\[
\det A = \sum_{\tau}\varepsilon(\tau)\prod_j a_{j,\tau(j)}
= \det(A^{\mathsf T}),
\]
末一和式即对转置矩阵的元素应用置换公式，
其中 $(A^{\mathsf T})_{ij} = a_{ji}$。

(4) 固定第 $j$ 列，由线性性把 $x_j = \sum_i a_{ij} e_i$ 拆开，
得 $\det A = \sum_i a_{ij}\, \det(\dots, e_i, \dots)$，
将 $e_i$ 移到最后一位（行的对换 $n - i$ 次，
列的对换 $n - j$ 次，由 (3) 得 $\det(\dots, e_i, \dots) =
(-1)^{i+j}\Delta_{ij}$，与子式的关系式正是
本科一年级分册的余子式法则。相似性：$\det(P^{-1}AP) = \det P^{-1}\det A \det P =
\det A$，由 (1) 可得。
\end{proof}

\begin{example}[余子式展开的具体计算]\label{ex:b2:linalg:cofactorwork}
计算
\[
\det\begin{pmatrix}
2 & 1 & 3\\
0 & 4 & 1\\
1 & 2 & 0
\end{pmatrix}
\]
沿第一列展开（此列只有一个零，偷懒的程度有限）。
符号按棋盘格 $(-1)^{i+j}$ 决定：
\[
2\,\det\begin{pmatrix}4 & 1\\ 2 & 0\end{pmatrix}
- 0
+ 1\cdot\det\begin{pmatrix}1 & 3\\ 4 & 1\end{pmatrix}
= 2(0 - 2) + (1 - 12) = -15 .
\]
用萨吕斯法则验算（见\cref{ex:b2:linalg:sarrus}）：$0 + 1 + 0
- 12 - 0 - 4 = -15$。策略而非教条：沿
含零最多的一行或一列展开；若都没有零，则
先用行变换造出一些零——一轮消元的代价
小于两层余子式展开。
\end{example}

\begin{example}[用运算法则计算行列式]\label{ex:b2:linalg:onesmatrix}
设 $J \in \mathcal{M}_n(K)$ 为元素全为 1 的矩阵，$a \in K$；
用刚证得的工具计算 $\det(aI_n + J)$。
$aI_n + J$ 的各列都以相同方式处理：把所有行加到第一行上
（行列式不变——把某一行的倍数加到另一行上
相当于加上一个重复方向的项，被交错性消去）。
第一行变为 $(a + n, a + n, \dots, a + n)$；由该行的线性性提出
公因子 $a + n$，再从其余每一列中减去第一列，
所得矩阵是三角矩阵，对角元为
$(1, a, \dots, a)$。于是
\[
\det(aI_n + J) = (a + n)\,a^{\,n-1}.
\]
收尾的观察：根 $a = 0$（重数为 $n - 1$）和
$a = -n$ 说明 $J$ 有特征值 $0$（重数 $n - 1$）
及特征值 $n$（重数为一）——秩为一矩阵
$J$ 的谱，提前一章（\cref{ch:b2:reduction}
将使之系统化）。
\end{example}

\begin{example}[用置换公式计算行列式]\label{ex:b2:linalg:permexample}
对于零元素较多的矩阵，此公式本身就很实用。
例如，在
\[
A = \begin{pmatrix}
0 & a & 0 & 0\\
0 & 0 & b & 0\\
0 & 0 & 0 & c\\
d & 0 & 0 & 0
\end{pmatrix},
\]
唯一选出非零元素的置换是 $4$-轮换
$\sigma = (1\,2\,3\,4)$，将列 $1 \to$ 行 $4$，等等；
$\varepsilon(\sigma) = (-1)^3 = -1$，故 $\det A = -abcd$。（验算
可经三次列对换化为对角矩阵。）
\end{example}

\begin{example}[用乘积公式计算范德蒙行列式]\label{ex:b2:linalg:vandermondework}
对于节点 $0, 1, 2$（求积公式如
\cref{exo:b2:linalg:4}所用的那样），由\cref{exo:b2:linalg:11}中的
范德蒙行列式一眼即可算出：
\[
\det\begin{pmatrix}
1 & 1 & 1\\
0 & 1 & 2\\
0 & 1 & 4
\end{pmatrix}
= (1 - 0)(2 - 0)(2 - 1) = 2 ,
\]
按第一列直接展开得 $1\cdot(4 - 2)
= 2$，两者一致。互异节点处行列式非零，即是
插值理论全部内容在一个行列式中的体现：
求值型 $P \mapsto P(a_i)$ 构成对偶空间的一组基，当且仅当
该行列式非零，即对互异的 $a_i$ 总成立
--- 这正是\cref{ex:b2:linalg:dualexamples}的定量化。
\end{example}

\section{迹，再论}

\begin{proposition}\label{prop:b2:linalg:trace}
迹 $\operatorname{tr} \colon \mathcal{M}_n(K) \to K$ 是满足
$\operatorname{tr}(AB) =
\operatorname{tr}(BA)$ 与 $\operatorname{tr}(I_n) = n$ 的唯一线性型（当
$\operatorname{char} K = 0$ 时）；自同态的迹
借助任一矩阵表示是良定的，且
\[
\operatorname{tr}(u) = \sum_{i} e_i^*\bigl(u(e_i)\bigr)
\]
在任一基下 --- 对偶性以不依赖基的方式写出迹。
\end{proposition}

\begin{proof}
$\operatorname{tr}(AB) = \operatorname{tr}(BA)$ 及基不变性
已在本科一年级证明。唯一性：线性型 $t$ 满足 $t(AB) =
t(BA)$，它零化每个换位子 $AB - BA$。我们断言换位子
张成迹为零的超平面，其维数为 $n^2 - 1$。只需两组
换位子即可。基本矩阵的乘法法则为
$E_{ab}E_{cd} = \delta_{bc}E_{ad}$。对
$i \neq j$ 时，得
\[
E_{ii}E_{ij} - E_{ij}E_{ii} = E_{ij} - 0 = E_{ij}
\]
（第二项 $E_{ij}E_{ii} = \delta_{ji}E_{ii} = 0$
由 $j \neq i$ 可知）：每个非对角的 $E_{ij}$ 都是换位子。
又有
\[
E_{ij}E_{ji} - E_{ji}E_{ij} = E_{ii} - E_{jj} .
\]
$E_{ij}$（$i \neq j$，共 $n^2 - n$ 个）连同
$E_{11} - E_{jj}$（$j \geq 2$，共 $n - 1$ 个）构成
$n^2 - 1$ 个线性无关的迹为零的矩阵，它们张成
超平面 $\ker\operatorname{tr}$。故 $t$ 在
$\operatorname{tr}$ 为零处亦为零，并经由它分解：$t =
c\operatorname{tr}$；再由 $t(I) = n$ 得 $c = 1$。下式：
第 $i$ 个对角元（取 $u$ 的矩阵）恰为
$e_i^*(u(e_i))$.
\end{proof}

\begin{remark}[常见误区]\label{rem:b2:linalg:pitfalls}
（i）行列式关于\emph{列}是 $n$ 重线性的，而非
关于矩阵线性：一般而言 $\det(A + B) \neq \det A + \det B$，
并且 $\det(\lambda A) = \lambda^n\det A$，而非
$\lambda\det A$。（ii）转置颠倒乘积的次序：
$(vu)^{\mathsf T} = u^{\mathsf T}v^{\mathsf T}$；忘记
乘积次序的颠倒会毁掉所有涉及逆的计算。（iii）零化子 $F^\circ$
位于 $E^*$ 中而非 $E$ 中，它
只有在内积之后，才成为熟悉的“正交补”，
而内积使 $E$ 与 $E^*$ 等同（见\cref{ch:b2:quadratic}）；这样的
等同并非典范的。（iv）“行秩等于列秩
”并不意味着行\emph{空间}与列空间相等 ---
二者位于不同的空间（$K^n$ 与 $K^m$）中，并
通过\cref{prop:b2:linalg:transposerank}相联系，而不相等。
（v）置换公式只是证明的工具：对于具体的数，应
用行变换与代数余子式（\cref{ex:b2:linalg:sarrus}）。
\end{remark}

\begin{example}[迹配对分解矩阵空间]\label{ex:b2:linalg:tracepairing}
在 $\mathcal{M}_2(\R)$ 上，取\cref{exo:b2:linalg:9}中的配对 $\langle A, B\rangle =
\operatorname{tr}(AB)$，将 $M =
\begin{psmallmatrix}1 & 4\\ 2 & 3\end{psmallmatrix}$ 分解为
对称部分与反对称部分，
\[
M = S + A, \qquad
S = \tfrac12(M + M^{\mathsf T}) =
\begin{pmatrix}1 & 3\\ 3 & 3\end{pmatrix},
\qquad
A = \tfrac12(M - M^{\mathsf T}) =
\begin{pmatrix}0 & 1\\ -1 & 0\end{pmatrix}.
\]
于是 $\operatorname{tr}(SA) = \operatorname{tr}
\begin{psmallmatrix}-3 & 1\\ -3 & 3\end{psmallmatrix} = 0$，两部分
关于迹配对是“正交”的 --- 这是一般事实的一个实例：
该事实在本章周末问题中证明，即反对称矩阵恰好构成
对称矩阵的零化子。
对偶性在选取任何内积之前，即已揭示
分解 $\mathcal{M}_n = \mathcal{S}_n \oplus
\mathcal{A}_n$。
\end{example}

\begin{remark}[本卷中的几种视角]\label{rem:b2:linalg:perspectives}
留意本章三个构造在后文中如何
变换形式。\emph{转置}将重新出现于
\cref{ch:b2:reduction}中：$u$ 与 $u^{\mathsf T}$ 具有
相同的特征值，且几何重数相等（本章的
周末问题第 15 题），这正是矩阵的行分析与列
分析永不相互矛盾。\emph{行列式}
在\cref{ch:b2:reduction}中化为关于参数的函数
（$\chi_u(X) = \det(X\,\mathrm{id} - u)$），并在
\cref{ch:b2:multint}中成为雅可比行列式，其多重线性在此化为
换元因子。\emph{迹}孕育了相似不变量：
它是 $\chi_u$ 的第二个系数，
即特征值之和，并最终
化为\cref{ch:b2:fourier}风格恒等式中对角线的积分。
一章线性代数，三道长长的影子。
\end{remark}

\begin{remark}[本章的应用]
对偶空间并非为抽象而抽象：零化子与转置
支撑着线性方程组的可解性理论（本章的周末问题
正是由它们证明了有限维情形的弗雷德霍姆择一定理），非退化配对
在\cref{ch:b2:quadratic}中以极化形式重新出现，
在\cref{ch:b2:hermitian}中化为伴随，
而这里构造的行列式支撑着\cref{ch:b2:reduction}的全部内容。
在本科三年级分册中，同一对偶性被推广到无穷维，
化为希尔伯特空间上的里斯表示定理与
弗雷德霍姆理论，此时以紧性
取代这里所用的
维数计数。
\end{remark}

\section{练习}

\begin{exercise}[$\star$]\label{exo:b2:linalg:1}
在 $\R^3$ 中，设 $\varphi_1(x,y,z) = x + y$，$\varphi_2 = y + z$，
$\varphi_3 = x + z$。证明 $(\varphi_1, \varphi_2, \varphi_3)$
是 $(\R^3)^*$ 的一组基，并求 $\R^3$ 的一组基，
使上述基为其对偶基。
\end{exercise}

\begin{exercise}[$\star$]\label{exo:b2:linalg:2}
用排列公式计算下列矩阵的行列式，
\[
\begin{pmatrix} 0 & 0 & a\\ 0 & b & 0\\ c & 0 & 0 \end{pmatrix},
\qquad
\begin{pmatrix}
a & b & 0 & 0\\
c & d & 0 & 0\\
0 & 0 & e & f\\
0 & 0 & g & h
\end{pmatrix},
\]
并写出第二个矩阵所提示的分块对角规则。
\end{exercise}

\begin{exercise}[$\star$]\label{exo:b2:linalg:3}
设 $F = \{(x,y,z,t) \in \R^4 : x + y = z + t \text{ 且 } x = 2y\}$。
求 $F^\circ$ 的一组基，并
在维数上检验 \cref{thm:b2:linalg:annihilator}。
\end{exercise}

\begin{exercise}[$\star\star$]\label{exo:b2:linalg:4}
设 $a_0, \dots, a_n$ 为 $K$ 中互不相同的点，且令 $\varphi_i \colon
P \mapsto P(a_i)$ 为 $K_n[X]$ 上的线性形式。证明 $(\varphi_0, \dots,
\varphi_n)$ 是 $K_n[X]^*$ 的一组基，并求出以之为对偶基的基，
再将形式 $P \mapsto \int_0^1 P(t)\,\dd t$（取 $K = \R$，
$n = 2$，$a_i = 0, \frac12, 1$）在该基下展开，
并由此识别出辛普森公式。
\end{exercise}

\begin{exercise}[$\star\star$]\label{exo:b2:linalg:5}
设 $u \in \mathcal{L}(E)$，$\dim E = n$，且 $\operatorname{rk} u
= 1$。证明 $u = \varphi(\cdot)\, a$，其中 $a$ 为某个向量，
$\varphi$ 为某个线性形式；再证 $\operatorname{tr} u = \varphi(a)$，
且 $u^2 = (\operatorname{tr} u)\, u$。由此推出 $\det(I + u) = 1 +
\operatorname{tr} u$。
\end{exercise}

\begin{exercise}[$\star\star$]\label{exo:b2:linalg:6}
证明 $\mathcal{M}_n(K)$（$n \geq 2$）的每个超平面
都包含一个可逆矩阵。\emph{提示：超平面为 $\{M :
\operatorname{tr}(AM) = 0\}$，对某个 $A \neq 0$
（见\cref{exo:b2:linalg:9}）。若 $A$ 为数量矩阵，则可取一个迹为零的可逆矩阵；
否则求可逆矩阵 $M$，使
$AM$ 的对角元全为零——类置换矩阵即可。}
\end{exercise}

\begin{exercise}[$\star\star$]\label{exo:b2:linalg:7}
（行列式的导数）设 $A \in \mathcal{M}_n(\R)$，
利用多重线性证明下式：
\[
\frac{\dd}{\dd t}\Big|_{t=0} \det(I_n + tA) = \operatorname{tr} A ,
\]
并推出 $\det(\eu^{tA}) = \eu^{t\operatorname{tr} A}$，假设
$t \mapsto \det(\eu^{tA})$ 可微，且群性质
$\eu^{(s+t)A} = \eu^{sA}\eu^{tA}$（已在
\cref{ch:b2:diffeq}中建立）。
\end{exercise}

\begin{exercise}[$\star\star$]\label{exo:b2:linalg:8}
（循环矩阵，$3 \times 3$）设 $j = \eu^{2\iu\pi/3}$，并令
\[
C = \begin{pmatrix}
a & b & c\\
c & a & b\\
b & c & a
\end{pmatrix} \in \mathcal{M}_3(\C).
\]
验证 $1, j, j^2$ 的范德蒙德矩阵的各列是
$C$ 的特征向量，并由此推出
\[
\det C = (a + b + c)(a + bj + cj^2)(a + bj^2 + cj).
\]
\end{exercise}

\begin{exercise}[$\star\star\star$]\label{exo:b2:linalg:9}
证明每个线性形式 $t$（定义在 $\mathcal{M}_n(K)$ 上）都可写成 $M \mapsto
\operatorname{tr}(AM)$ 的形式，其中 $A$ 唯一；映射 $A \mapsto
\operatorname{tr}(A\,\cdot)$ 是从
$\mathcal{M}_n(K)$ 到其对偶空间的同构。由此再次推出
\cref{prop:b2:linalg:trace}的唯一性结论。
\end{exercise}

\begin{exercise}[$\star\star\star$]\label{exo:b2:linalg:10}
设 $u, v \in \mathcal{L}(E)$，满足 $u \circ v - v \circ u = u$。
证明 $u$ 幂零。\emph{提示：先证明
$\operatorname{tr}(u^k) = 0$（对所有 $k \geq 1$，用归纳法计算 $u^k v -
v u^k$），再利用下面的事实，其证明可借助
牛顿恒等式或对维数归纳：
$\C$ 上向量空间的自同态，若其所有幂的迹均为零，则它幂零。
在 $\C$ 上讨论。}
\end{exercise}

\begin{exercise}[$\star\star$]\label{exo:b2:linalg:11}
（范德蒙德）对 $a_0, \dots, a_n \in K$，证明下式：
\[
\det\begin{pmatrix}
1 & 1 & \cdots & 1\\
a_0 & a_1 & \cdots & a_n\\
\vdots & \vdots & & \vdots\\
a_0^n & a_1^n & \cdots & a_n^n
\end{pmatrix}
= \prod_{0 \leq i < j \leq n} (a_j - a_i).
\]
\emph{（将行列式视为 $a_n$ 的多项式：确定其
次数、根与首项系数；再用归纳法。）}
\end{exercise}

\begin{exercise}[$\star\star\star$]\label{exo:b2:linalg:12}
设 $A, B, C, D \in \mathcal{M}_n(K)$，$K$ 为无限域，且
满足 $CD = DC$。证明下式成立：
\[
\det\begin{pmatrix} A & B\\ C & D\end{pmatrix}
= \det(AD - BC).
\]
\emph{（先设 $D$ 可逆，在右侧乘以
$\begin{psmallmatrix} I & 0\\ -D^{-1}C & I\end{psmallmatrix}$;
再将 $D$ 换为 $D + tI$，比较关于 $t$ 的两个多项式。）}
\end{exercise}

\section{问题：弗雷德霍姆择一性}

线性方程组 $u(x) = b$ 何时有解？
完整的回答是一个对偶性陈述：\emph{当且仅当 $b$ 被
所有零化其像的线性形式零化时
$u$} --- 这些形式可以计算，它们恰为
转置的核。本周末问题建立有限维对偶的完整
“词典”（包括形式的分解、双对偶、
零化子演算、转置），证明有限维
\emph{弗雷德霍姆择一性}，并以
迹形式与一个刻画作结：迹是相似下唯一的
线性不变量。下文中，$E$ 与 $F$ 均为
$K$ 上的有限维向量空间，$n = \dim E$。

\begin{problem}[{周末问题 --- 有限维中的对偶
与弗雷德霍姆择一性}]
\label{pb:b2:linalg:1}
记号：对 $S \subseteq E^*$，\emph{前零化子}为
$S_\circ = \{x \in E : \varphi(x) = 0 \text{ 对所有 } \varphi
\in S\}$；零化子 $F^\circ$ 与转置 $u^{\mathsf T}$
分别见 \cref{def:b2:linalg:annihilator} 与
\cref{def:b2:linalg:transpose}。

\medskip
\noindent\textbf{第一部分 --- 分解引理。}设
$\varphi_1, \dots, \varphi_p, \varphi \in E^*$。
\begin{enumerate}
  \item 设 $\Phi \colon E \to K^p$，$x \mapsto (\varphi_1(x),
        \dots, \varphi_p(x))$。确定 $\ker\Phi$，并证明
        $\Phi^{\mathsf T}$ 将 $K^p$ 的坐标形式映为
        $\varphi_i$，并推得
        \[
        \dim \bigl(\ker\varphi_1 \cap \dots \cap
        \ker\varphi_p\bigr) = n - \dim
        \operatorname{Vect}(\varphi_1, \dots, \varphi_p).
        \]
  \item （分解引理）证明下述等价：
        \[
        \varphi \in \operatorname{Vect}(\varphi_1, \dots,
        \varphi_p)
        \iff
        \ker\varphi_1 \cap \dots \cap \ker\varphi_p \subseteq
        \ker\varphi .
        \]
  \item 由此推出：$(\varphi_1, \dots, \varphi_p)$ 线性无关当且仅当
        $\bigcap_i \ker\varphi_i$ 的维数为 $n - p$；且
        余维数为 $p$ 的子空间是 $p$ 个
        超平面的交，不能更少。
  \item 在 $\R^4$ 中，设 $\varphi_1 = x + y - z$，$\varphi_2 = y +
        z - t$，$\psi = x + 2y - t$ 与 $\psi' = x + y + t$。
        由分解引理，判断 $\psi$ 与
        $\psi'$ 是否属于 $\operatorname{Vect}(\varphi_1,
        \varphi_2)$。
  \item 在 $E = \R_2[X]$ 上，证明 $\psi_0 \colon P \mapsto
        P(0)$，$\psi_1 \colon P \mapsto P(1)$，$\psi_2 \colon P
        \mapsto \int_0^1 P(t)\dd t$ 构成 $E^*$ 的一组基，
        求 $(P_0, P_1, P_2)$ 为 $E$ 的一组基，使上述泛函为其对偶基，
        并求唯一的 $P \in \R_2[X]$，满足 $P(0)
        = 1$，$P(1) = 2$，$\int_0^1 P = \frac32$。
\end{enumerate}

\medskip
\noindent\textbf{第二部分 --- 双对偶与零化子
演算。}
\begin{enumerate}[resume]
  \item 证明\emph{求值映射} $J \colon E \to
        E^{**}$，$J(x)(\varphi) = \varphi(x)$，是线性的且
        单射，从而在有限维下为同构。
  \item （双零化子）证明 $J(F) = F^{\circ\circ} :=
        (F^\circ)^\circ$ 对每个子空间 $F \subseteq E$ 成立：
        经 $J$ 等同后，零化子的零化子
        恰为该子空间本身。
  \item 证明零化子演算：$(F + G)^\circ = F^\circ
        \cap G^\circ$ 与 $(F \cap G)^\circ = F^\circ + G^\circ$。
  \item 推出（并直接重新证明）：两个核相同的非零
        线性形式彼此成比例。
  \item （前对偶基）证明：对 $(\varphi_1,
        \dots, \varphi_n)$ 为 $E^*$ 的任一基，存在唯一的
        $(u_1, \dots, u_n)$ 为 $E$ 的基，使 $\varphi_i(u_j) =
        \delta_{ij}$。
\end{enumerate}

\medskip
\noindent\textbf{第三部分 --- 转置演算。}
\begin{enumerate}[resume]
  \item 证明 $u \mapsto u^{\mathsf T}$ 是线性双射，
        从 $\mathcal{L}(E, F)$ 到 $\mathcal{L}(F^*, E^*)$ 上的，
        并证明 $(u^{-1})^{\mathsf T} = (u^{\mathsf T})^{-1}$，
        当 $u$ 可逆时。
  \item （自然性）证明 $u^{\mathsf T\mathsf T} \circ J_E
        = J_F \circ u$：在求值同构下，
        二次转置\emph{恰为} $u$。
  \item 证明：$u$ 满射当且仅当 $u^{\mathsf T}$ 单射；
        $u$ 单射当且仅当 $u^{\mathsf T}$ 满射。
  \item 对 $u \in \mathcal{L}(E)$，子空间 $F$ 在
        $u$ 下稳定，当且仅当 $F^\circ$ 在
        $u^{\mathsf T}$ 下稳定。
  \item 证明 $\ker(u^{\mathsf T} - \lambda\,
        \mathrm{id}_{E^*}) = \bigl(\operatorname{im}(u - \lambda\,
        \mathrm{id}_E)\bigr)^\circ$，并由此推出 $u$ 与
        $u^{\mathsf T}$ 的特征值相同，且
        几何重数相同。
\end{enumerate}

\medskip
\noindent\textbf{第 IV 部分 --- 弗雷德霍姆择一性。}
\begin{enumerate}[resume]
  \item 证明 $\operatorname{im} u = (\ker u^{\mathsf
        T})_\circ$，其中 $u \in \mathcal{L}(E, F)$，并由此推出有限维情形下的
        \emph{弗雷德霍姆择一性}：方程
        $u(x) = b$ 有解，当且仅当对每个 $\psi \in F^*$
        （满足 $u^{\mathsf T}\psi = 0$）都有
        $\psi(b) = 0$。
  \item 矩阵形式：对 $A \in \mathcal{M}_{m,n}(K)$ 与 $b \in
        K^m$，以下两者恰有一个成立：（i）$Ax = b$
        有解；（ii）存在 $y \in K^m$，使得
        $A^{\mathsf T}y = 0$ 且 $y^{\mathsf T}b = 1$。证明其中的
        “至多一个”与“至少一个”。
  \item 求所有使下面的方程组有解的 $b \in \R^3$，
        \[
        x + y = b_1, \qquad y + z = b_2, \qquad x + 2y + z = b_3
        \]
        通过计算转置
        矩阵的核来求解。
  \item （离散诺伊曼问题）在 $E = \R^n$（$n \geq 3$）上，
        定义 $L$，使 $(Lx)_k = x_k - \frac12(x_{k-1} +
        x_{k+1})$，下标按模 $n$ 理解。证明 $L^{\mathsf T} = L$
        （取典范同一化），证明 $\ker L$ 是
        由常向量组成的直线\emph{（考虑最大分量）}，
        并得出结论：$Lx = b$ 可解当且仅当 $\sum_k b_k = 0$。
\end{enumerate}

\medskip
\noindent\textbf{第 V 部分 --- 迹形式与
不变性定理。} 由\cref{exo:b2:linalg:9}可知，映射 $A \mapsto
\operatorname{tr}(A\,\cdot)$ 将 $\mathcal{M}_n(K)$ 与其
        对偶空间等同。假设 $\operatorname{char} K = 0$（例如 $K = \Q, \R,
\C$）。
\begin{enumerate}[resume]
  \item 在此同一化下，证明
        对称矩阵子空间 $\mathcal{S}_n$ 的零化子
        是反对称矩阵子空间 $\mathcal{A}_n$，
        且反之亦然。
  \item 证明超平面的零化子
        $\mathfrak{sl}_n = \{M : \operatorname{tr} M = 0\}$ 是
        直线 $K I_n$；等价地，在所有迹为零的矩阵上取零值的
        线性型是迹的倍数。
  \item 证明 $\mathcal{M}_n(K)$ 中每个矩阵都
        可写成两个可逆矩阵之和。
  \item （迹是唯一的线性相似不变量）设
        $t$ 为 $\mathcal{M}_n(K)$ 上的线性型，满足
        $t(PMP^{-1}) = t(M)$ 对所有 $M$ 及每个可逆矩阵
        $P$ 成立。先证 $t(PX) = t(XP)$（$P$ 可逆），然后
        证明 $t(BX) = t(XB)$ 对\emph{所有} $B$ 成立，并得出 $t = c
        \operatorname{tr}$，其中 $c \in K$。
  \item 证明 $\operatorname{rk} u \leq r$ 当且仅当 $u$
        是 $r$ 个秩 $\leq 1$ 的映射之和，即 $u =
        \sum_{i=1}^{r} \psi_i(\cdot)\,f_i$，其中 $\psi_i \in E^*$，
        $f_i \in F$；由此推出 $\operatorname{rk}(u + v) \leq
        \operatorname{rk} u + \operatorname{rk} v$。
  \item （综合）整理本题所证明的“字典”：
        子空间与零化子，和与
        交，映射与转置，可解性与
        转置核的正交性，迹与
        相似。对表中每一项，指出证明它的小题，
        并用一句话说明：维数变为无穷时，何物取代
        维数的计数（本科三年级分册在希尔伯特空间上
        将此讲得更为精确）。
\end{enumerate}
\end{problem}
